\section{The set of real numbers}

This section concerns the rules of arithmetic, inequalities, and absolute
values. We work with the usual operations on real numbers. In particular,
\emph{an inverse must belong to the set under discussion}: an integer does
not have a multiplicative inverse \emph{in the integers} merely because it
has one in the real numbers.

\subsection*{A reference sheet for the rules}
Ross's labels will be useful when explaining individual steps. For real
numbers $x,y,z$, the field rules are
\begin{align*}
 \mathrm{A1}:&\quad x+(y+z)=(x+y)+z,
 &\mathrm{A2}:&\quad x+y=y+x,\\
 \mathrm{A3}:&\quad x+0=x,
 &\mathrm{A4}:&\quad x+(-x)=0,\\
 \mathrm{M1}:&\quad x(yz)=(xy)z,
 &\mathrm{M2}:&\quad xy=yx,\\
 \mathrm{M3}:&\quad x\cdot1=x,
 &\mathrm{M4}:&\quad xx^{-1}=1\quad(x\ne0),\\
 \mathrm{DL}:&\quad x(y+z)=xy+xz.
\end{align*}
The additive identity $0$, the multiplicative identity $1$, the additive
inverse $-x$, and (when $x\ne0$) the multiplicative inverse $x^{-1}$ must
be elements of the field. A field has more than one element.

The order rules say that any two real numbers can be compared (O1),
$x\le y$ and $y\le x$ together imply $x=y$ (O2), and
$x\le y\le z$ implies $x\le z$ (O3). They also give
\[
 \mathrm{O4}:\quad x\le y\ \Longrightarrow\ x+z\le y+z,
 \qquad
 \mathrm{O5}:\quad x\le y,\ 0\le z\ \Longrightarrow\ xz\le yz.
\]
By definition, $x<y$ means that $x\le y$ and $x\ne y$.
Subtraction is addition of the additive inverse:
$x-y=x+(-y)$. Division by a nonzero number is multiplication by its
multiplicative inverse: $x/y=xy^{-1}$.

We may use the results already established in the text before these
exercises, including additive cancellation, $x\cdot0=0$,
$(-x)y=-(xy)$, nonnegativity of squares, and the triangle inequality.
When an exercise asks for a missing part of a theorem, that missing part
is proved below rather than assumed.

\exercise{3.1}{Which rules fail for natural numbers and integers?}

\goal Determine which of A1--A4, M1--M4, DL, and O1--O5 are valid when the
operations and order are the usual ones and all required elements must lie
in the specified set.

\step{(a), Step 1: fix the convention for natural numbers.}
Ross uses $\N=\{1,2,3,\ldots\}$. Thus $1\in\N$ but $0\notin\N$.
This distinction matters for the identity requirement A3.

\step{(a), Step 2: test the additive identity.}
A3 requires an additive identity \emph{in the set}. The usual identity is
$0$, which is not in $\N$. There is no replacement positive integer $e$:
if $e\in\N$, then $e\ge1$, so
\[
 1+e\ge1+1=2>1.
\]
Consequently $1+e\ne1$. Hence no natural number can act as an additive
identity, and the internal identity requirement A3 fails.

There is a small interpretive issue here. The equality $n+0=n$ is true if
we calculate in $\Z$ or $\R$. It is not a false numerical equation.
It fails as a \emph{field axiom for $\N$} because its required constant
$0$ is outside $\N$. We use this internal interpretation in the answer.

\step{(a), Step 3: test additive inverses.}
A4 would require an element of $\N$ that can be added to $1$ to give $0$.
But if $m\in\N$, then $m\ge1$, and therefore $1+m\ge2$.
In particular, $1+m\ne0$. The additive inverse $-1$ exists in $\Z$ and
$\R$, but not in $\N$. Thus A4 fails.

\step{(a), Step 4: test multiplicative inverses.}
M4 would require an element $m\in\N$ with $2m=1$.
Every such $m$ satisfies $m\ge1$, so $2m\ge2$. Therefore $2m\ne1$.
Equivalently, the required inverse $1/2$ is not a natural number.
Thus M4 fails. It is enough to find the single nonzero element $2$ for
which the requirement fails; M4 must hold for \emph{every} nonzero element.

\step{(a), Step 5: check the remaining rules.}
Addition and multiplication of positive integers again give positive
integers. Associativity and commutativity remain true, so A1, A2, M1,
and M2 hold. Since $1\in\N$ and $n\cdot1=n$, M3 holds. The distributive
identity holds for positive integers, so DL holds. The usual comparisons
of natural numbers satisfy O1--O4. Finally, multiplying an inequality
between natural numbers by a natural number, which is positive, preserves
its direction; this verifies the numerical inequality in O5.
The occurrence of $0$ in the condition of O5 is understood in the ambient
usual order. No additional order inequality becomes false upon
restricting its variables to $\N$.

\answer{For $\N$ the failed internal field requirements are
\(\boxed{\mathrm{A3},\ \mathrm{A4},\ \mathrm{M4}}\).
The remaining listed arithmetic identities and order implications are
valid with the usual arithmetic and order.}

\step{(b), Step 1: identify what changes for integers.}
Now the set is $\Z$, so $0\in\Z$ and $-n\in\Z$ whenever $n\in\Z$.
Thus $n+0=n$ takes place inside the set, and $n+(-n)=0$ also takes place
inside the set. A3 and A4 both hold.

\step{(b), Step 2: exhibit the remaining failure.}
The element $2$ is a nonzero integer. If it had an inverse $m\in\Z$,
we would have $2m=1$. The left side is an even integer, since it is twice
an integer; the right side is odd. Hence no such integer $m$ exists.
This again says that $1/2\notin\Z$, so M4 fails.

\step{(b), Step 3: check every other requirement.}
The sum and product of integers are integers. The associative,
commutative, and distributive identities remain valid. Both identities,
$0$ and $1$, are integers. Additive inverses are integers. The usual
order of integers satisfies O1--O4, and multiplication by a nonnegative
integer preserves a weak inequality, as required by O5.

\answer{For $\Z$ the only failed listed requirement is
\(\boxed{\mathrm{M4}}\). Some integers, namely $1$ and $-1$, do have
integer multiplicative inverses, but that does not meet a requirement
imposed on every nonzero integer.}

\exercise{3.2}{Locate the hidden uses of additive commutativity}

\goal Find the uses of A2 in the proofs of Theorem~3.1(ii) and (iii).

\idea The cancellation rule proved in Theorem~3.1(i) is written with the
common term on the \emph{right}:
\[
 x+c=y+c\quad\Longrightarrow\quad x=y.
\]
Before using this particular form, both equalities must have that shape.

\step{(a), Step 1: write what distributivity actually gives.}
Set $t=a\cdot0$. Since $0+0=0$ by A3, distributivity gives
\[
 t=a\cdot0=a(0+0)=a\cdot0+a\cdot0=t+t.
\]
To use right cancellation with the common term $t$, the proof rewrites
the leftmost $t$ as $0+t$.

\step{(a), Step 2: justify that rewriting.}
A3 gives $t+0=t$, with the zero on the \emph{right}. A2 is needed to move
it to the left:
\[
 0+t=t+0=t.
\]
Thus $t=t+t$ becomes $0+t=t+t$. Right cancellation of $t$ now gives
$0=t=a\cdot0$.

\answer{The hidden use of A2 is the passage from $a\cdot0$ to
$0+a\cdot0$: first A3 permits $a\cdot0+0$, then A2 switches the two
summands.}

\step{(b), Step 1: identify the equality just before cancellation.}
The proof obtains
\[
 ab+(-a)b=ab+\bigl(-(ab)\bigr).
\]
Here the common summand $ab$ is on the \emph{left}, whereas the previously
proved cancellation statement has the common summand on the right.

\step{(b), Step 2: apply A2 to both sums.}
The two relevant switches are
\[
 ab+(-a)b=(-a)b+ab,
 \qquad
 ab+\bigl(-(ab)\bigr)=\bigl(-(ab)\bigr)+ab.
\]
The equality therefore becomes
\[
 (-a)b+ab=\bigl(-(ab)\bigr)+ab.
\]
Now right cancellation with common term $ab$ proves
$(-a)b=-(ab)$.

\step{(b), Step 3: distinguish addition from multiplication.}
There are also commutative switches in the preceding factorization, but
those use M2, not A2. For example,
\begin{align*}
 ab+(-a)b
   &=ba+b(-a) &&\text{by M2 in the two products},\\
   &=b\bigl(a+(-a)\bigr) &&\text{by DL, read in reverse},\\
   &=b\cdot0 &&\text{by A4},\\
   &=0 &&\text{by Theorem~3.1(ii)}.
\end{align*}
No exchange of \emph{summands} occurs in this displayed chain.

\answer{The A2 needed for the final use of Theorem~3.1(i) moves the
common additive term $ab$ from the left to the right on both sides.
Product switches elsewhere in the argument use M2.}

\exercise{3.3}{A product of two negatives and multiplicative cancellation}

\goal Prove Theorem~3.1(iv), $(-a)(-b)=ab$, and Theorem~3.1(v),
$ac=bc$ with $c\ne0$ implies $a=b$.

\step{Part (iv), Step 1: prove that taking the negative twice returns the number.}
For any real $u$, A4 gives
\[
 u+(-u)=0.
\]
Apply A4 again, now to the number $-u$:
\[
 (-u)+\bigl(-(-u)\bigr)=0.
\]
By A2 the second left side is $-(-u)+(-u)$. Hence
\[
 u+(-u)=-(-u)+(-u).
\]
The common right summand is $-u$. Additive cancellation gives
\[
 u=-(-u).
\]
This proves the double-negative rule without using a product of two
negative numbers.

\step{Part (iv), Step 2: reduce the product to the rule already proved.}
Theorem~3.1(iii) states $(-x)y=-(xy)$ for all real $x,y$.
First apply it with $x=a$ and $y=-b$:
\[
 (-a)(-b)=-\bigl(a(-b)\bigr).
\]
Next rewrite the product inside the parentheses:
\begin{align*}
 a(-b)
   &=(-b)a &&\text{by M2},\\
   &=-(ba) &&\text{by Theorem~3.1(iii) with $x=b$, $y=a$},\\
   &=-(ab) &&\text{by M2 inside the negative}.
\end{align*}
Substitute this expression into the previous equality:
\[
 (-a)(-b)=-\bigl(-(ab)\bigr)=ab.
\]
The last equality is Step~1 with $u=ab$.

\answer{For all real $a,b$, \(\boxed{(-a)(-b)=ab}\).}

\step{Part (v), Step 1: use the nonzero assumption.}
Suppose $ac=bc$ and $c\ne0$. By M4 the number $c^{-1}$ exists and
satisfies $cc^{-1}=1$. Multiply both sides of the given equality on the
right by this \emph{same} real number:
\[
 (ac)c^{-1}=(bc)c^{-1}.
\]

\step{Part (v), Step 2: simplify each side, naming each rule.}
For the left side,
\[
 (ac)c^{-1}=a(cc^{-1})=a\cdot1=a.
\]
These equalities use M1, M4, and M3, respectively. The same rules give
\[
 (bc)c^{-1}=b(cc^{-1})=b\cdot1=b.
\]
Replacing the two sides of $(ac)c^{-1}=(bc)c^{-1}$ by these simplified
forms gives $a=b$.

\answer{If $ac=bc$ and $c\ne0$, then \(\boxed{a=b}\).
The condition $c\ne0$ is necessary: $1\cdot0=2\cdot0$, but $1\ne2$.}

\exercise{3.4}{Positivity of one and reversal of reciprocals}

\goal Prove Theorem~3.2(v), $0<1$, and Theorem~3.2(vii):
\[
 0<a<b\quad\Longrightarrow\quad0<b^{-1}<a^{-1}.
\]

\step{Part (v), Step 1: show that $1$ is nonnegative.}
Theorem~3.2(iv), already proved in the text, says $0\le x^2$ for every
real $x$. Take $x=1$. By M3, $1^2=1\cdot1=1$. Therefore
\[
 0\le1.
\]

\step{Part (v), Step 2: rule out equality with zero.}
Suppose, for a contradiction, that $1=0$. Let $x$ be any element of the
field. Then
\[
 x=x\cdot1=x\cdot0=0.
\]
The first equality uses M3; the second substitutes $1=0$; the third uses
Theorem~3.1(ii). Thus \emph{every} field element would equal $0$, leaving
only one element. This contradicts the requirement that a field have
more than one element. Consequently $1\ne0$.

\answer{Since $0\le1$ and $0\ne1$, the definition of strict inequality
gives \(\boxed{0<1}\).}

\step{Part (vii), Step 1: record why the inverses are positive.}
If $t>0$, then $t\ne0$, so M4 supplies $t^{-1}$ with $tt^{-1}=1$.
Here is the proof of its positivity, as in Theorem~3.2(vi).
First $t^{-1}$ cannot be zero, because then
$1=tt^{-1}=t\cdot0=0$, contradicting Part~(v).
If $t^{-1}<0$, order reversal under negation gives $0<-t^{-1}$.
The product of nonnegative numbers is nonnegative
(Theorem~3.2(iii)), so
\[
 0\le t(-t^{-1})=-(tt^{-1})=-1.
\]
The product identity in the middle follows from Theorem~3.1(iii) and
M2. Negating $0\le-1$ reverses its direction and gives $1\le0$.
Together with $0\le1$, O2 would give $1=0$, again contradicting
Part~(v). Thus $t^{-1}$ is neither zero nor negative, so $t^{-1}>0$.

In the hypothesis $0<a<b$, we have $a>0$ and $b>0$.
Applying the preceding argument to $t=a$ and $t=b$ gives
\[
 a^{-1}>0,\qquad b^{-1}>0.
\]

\step{Part (vii), Step 2: justify multiplication of strict inequalities.}
If $x<y$ and $c>0$, then O5 first gives the weak inequality $xc\le yc$.
Equality $xc=yc$ would imply $x=y$ by Exercise~3.3(v), since $c\ne0$.
But $x<y$ includes $x\ne y$. Therefore equality is impossible, and
\[
 x<y,\quad c>0\quad\Longrightarrow\quad xc<yc.
\]
This explains why strictness is retained in the following two steps.

\step{Part (vii), Step 3: multiply first by $a^{-1}$.}
Start from $a<b$ and use the positive multiplier $a^{-1}$:
\[
 aa^{-1}<ba^{-1}.
\]
M4 turns the left side into $1$, so
\[
 1<ba^{-1}.
\]

\step{Part (vii), Step 4: multiply next by $b^{-1}$.}
The multiplier $b^{-1}$ is also positive. Thus
\[
 1\cdot b^{-1}<(ba^{-1})b^{-1}.
\]
The left side equals $b^{-1}$ by M2 and M3. On the right,
\[
 (ba^{-1})b^{-1}
   =(a^{-1}b)b^{-1}
   =a^{-1}(bb^{-1})
   =a^{-1}\cdot1
   =a^{-1}.
\]
These steps use M2, M1, M4, and M3, in that order. Hence
$b^{-1}<a^{-1}$.

\answer{Combining this strict inequality with $b^{-1}>0$ gives
\(\boxed{0<b^{-1}<a^{-1}}\). Positivity of $a$ and $b$ is essential
to the choice of positive multipliers.}

\exercise{3.5}{Absolute-value bounds and the reverse triangle inequality}

\goal Prove (a) $\abs{b}\le a$ if and only if $-a\le b\le a$, and
(b) $\abs{\abs{a}-\abs{b}}\le\abs{a-b}$.

\step{Part (a), Step 1: unpack the definition of absolute value.}
For a real number $x$,
\[
 \abs{x}=\begin{cases}x,&x\ge0,\\-x,&x<0.\end{cases}
\]
In both cases $\abs{x}\ge0$. We also have
\[
 -\abs{x}\le x\le\abs{x}.
\]
To verify the second statement rather than assume it, consider the two
cases. If $x\ge0$, then $\abs{x}=x$, the upper inequality is equality,
and the lower inequality is $-x\le x$, since $-x\le0\le x$.
If $x<0$, then $\abs{x}=-x$, the lower inequality is equality
$-\abs{x}=x$, and the upper inequality follows from $x<0<-x$.

\step{Part (a), Step 2: prove the forward implication.}
Assume $\abs{b}\le a$. The bounds from Step~1 give $b\le\abs{b}$.
Transitivity then gives
\[
 b\le\abs{b}\le a,
\]
so $b\le a$. Negating the assumed inequality reverses its direction:
\[
 -a\le-\abs{b}.
\]
Since $-\abs{b}\le b$ by Step~1, transitivity gives $-a\le b$.
The two conclusions together are $-a\le b\le a$.

\step{Part (a), Step 3: prove the reverse implication.}
Assume both $-a\le b$ and $b\le a$.
If $b\ge0$, then $\abs{b}=b$, so the second assumption gives
$\abs{b}\le a$.
If $b<0$, negate the first assumption $-a\le b$ to obtain
$-b\le a$. In this case $\abs{b}=-b$, so again $\abs{b}\le a$.
These two cases exhaust all real $b$ by O1.

\step{Part (a), Step 4: check zero and negative bounds.}
If $a=0$, the left condition says $\abs{b}\le0$, which forces
$\abs{b}=0$ and hence $b=0$ by the definition. The right condition is
$0\le b\le0$, which forces $b=0$ by O2.
If $a<0$, the left condition is impossible because
$0\le\abs{b}$, whereas the right condition is impossible because
$-a>0>a$: no real $b$ can simultaneously be at least $-a$ and at most
$a$. Thus the equivalence includes these bounds as well.

\answer{For all real $a,b$, including $a\le0$,
\(\boxed{\abs{b}\le a\ \Longleftrightarrow\ -a\le b\le a}\).}

\step{Part (b), Step 1: obtain a one-sided bound from the triangle inequality.}
The equality $a=(a-b)+b$ follows by associating $a+((-b)+b)$ and
using $(-b)+b=0$. Apply the triangle inequality to the sum $(a-b)+b$:
\[
 \abs{a}=\abs{(a-b)+b}\le\abs{a-b}+\abs{b}.
\]
Add $-\abs{b}$ to both sides using O4. After associating and cancelling
$\abs{b}+(-\abs{b})$, the result is
\[
 \abs{a}-\abs{b}\le\abs{a-b}. \tag{1}
\]

\step{Part (b), Step 2: bound the negative of the same expression.}
Interchange the roles of $a$ and $b$ in the preceding argument:
\[
 \abs{b}-\abs{a}\le\abs{b-a}. \tag{2}
\]
We need to identify the right side. First, $b-a=-(a-b)$, as follows by
distributing the negative over $a+(-b)$ and using the double-negative
rule. Also $\abs{-x}=\abs{x}$ for any real $x$: when $x>0$, the two
absolute values are $-(-x)=x$ and $x$; when $x<0$, both are $-x$;
when $x=0$, both are $0$. Therefore
\[
 \abs{b-a}=\abs{-(a-b)}=\abs{a-b}.
\]
Substitute this into (2):
\[
 \abs{b}-\abs{a}\le\abs{a-b}.
\]
Negating both sides reverses the inequality, so
\[
 -\abs{a-b}\le\abs{a}-\abs{b}. \tag{3}
\]

\step{Part (b), Step 3: combine and apply part (a).}
Together (1) and (3) give the two-sided bound
\[
 -\abs{a-b}\le\abs{a}-\abs{b}\le\abs{a-b}.
\]
In part (a), take its bounded number to be
$\abs{a}-\abs{b}$ and its bound to be $\abs{a-b}$. That equivalence
turns this two-sided bound into
\[
 \abs{\abs{a}-\abs{b}}\le\abs{a-b}.
\]

\answer{The reverse triangle inequality is
\(\boxed{\abs{\abs{a}-\abs{b}}\le\abs{a-b}}\) for every real $a,b$.}

\exercise{3.6}{The triangle inequality for three and for finitely many terms}

\goal Prove (a) $\abs{a+b+c}\le\abs{a}+\abs{b}+\abs{c}$ and
(b) $\abs{a_1+\cdots+a_n}\le\abs{a_1}+\cdots+\abs{a_n}$
for every positive integer $n$ and every choice of real summands.

\step{Part (a), Step 1: group the sum into two terms.}
By A1, $a+b+c=(a+b)+c$. The two-term triangle inequality, applied to
the real numbers $a+b$ and $c$, gives
\[
 \abs{a+b+c}\le\abs{a+b}+\abs{c}. \tag{4}
\]

\step{Part (a), Step 2: apply the two-term inequality again.}
Apply it to $a$ and $b$:
\[
 \abs{a+b}\le\abs{a}+\abs{b}.
\]
Add $\abs{c}$ to both sides by O4:
\[
 \abs{a+b}+\abs{c}\le\abs{a}+\abs{b}+\abs{c}. \tag{5}
\]
By transitivity, combining (4) with (5) gives
\[
 \abs{a+b+c}\le\abs{a}+\abs{b}+\abs{c}.
\]

\answer{The three-term inequality follows from two applications of the
two-term inequality; no enumeration of the signs of $a,b,c$ is needed.}

\step{Part (b), Step 1: state exactly what is being proved by induction.}
Let $P_n$ be the assertion that
\[
 \abs{a_1+\cdots+a_n}\le\abs{a_1}+\cdots+\abs{a_n}
\]
for \emph{every} list of $n$ real numbers $a_1,\ldots,a_n$.
This quantifier is important: in the induction step we may choose any
list of $n+1$ numbers and apply $P_n$ to its first $n$ entries.

\step{Part (b), Step 2: verify the base case.}
For $n=1$, the claimed inequality is simply $\abs{a_1}\le\abs{a_1}$.
The two sides are equal, so $P_1$ holds for every real $a_1$.

\step{Part (b), Step 3: state the induction hypothesis and the next target.}
Fix $n\ge1$ and assume $P_n$. Let
$a_1,\ldots,a_n,a_{n+1}$ be arbitrary real numbers. The target is
\[
 \abs{a_1+\cdots+a_n+a_{n+1}}
   \le\abs{a_1}+\cdots+\abs{a_n}+\abs{a_{n+1}}.
\]
Set $s=a_1+\cdots+a_n$. The induction hypothesis gives
$\abs{s}\le\sum_{j=1}^{n}\abs{a_j}$.

\step{Part (b), Step 4: prove the induction step.}
The two-term triangle inequality and then the induction hypothesis give
\begin{align*}
 \abs{a_1+\cdots+a_n+a_{n+1}}
   &=\abs{s+a_{n+1}} &&\text{by the definition of $s$ and A1},\\
   &\le\abs{s}+\abs{a_{n+1}} &&\text{by the triangle inequality},\\
   &\le\sum_{j=1}^{n}\abs{a_j}+\abs{a_{n+1}}
       &&\text{by the induction hypothesis and O4},\\
   &=\sum_{j=1}^{n+1}\abs{a_j} &&\text{by the definition of the sum}.
\end{align*}
The list of $n+1$ real numbers was arbitrary, so $P_{n+1}$ follows.

\answer{Since $P_1$ holds and $P_n$ implies $P_{n+1}$ for every
$n\ge1$, mathematical induction proves the finite-sum triangle
inequality for every $n\in\N$.}

\exercise{3.7}{Strict absolute-value bounds and intervals centered at a point}

\goal Establish all three equivalences:
\begin{align*}
 \text{(a)}\quad&\abs{b}<a
    &&\Longleftrightarrow&&-a<b<a,\\
 \text{(b)}\quad&\abs{a-b}<c
    &&\Longleftrightarrow&&b-c<a<b+c,\\
 \text{(c)}\quad&\abs{a-b}\le c
    &&\Longleftrightarrow&&b-c\le a\le b+c.
\end{align*}

\step{Part (a), Step 1: prove the forward implication.}
Suppose $\abs{b}<a$. From Exercise~3.5(a), Step~1, we know
$b\le\abs{b}$ and $-\abs{b}\le b$.
Thus
\[
 b\le\abs{b}<a\quad\Longrightarrow\quad b<a.
\]
Negating $\abs{b}<a$ gives $-a<-\abs{b}$. This strict reversal is
justified by Theorem~3.2(i) for the weak inequality; equality of the
negatives would imply equality of the originals by taking their
negatives again, so equality is excluded. Combining it with
$-\abs{b}\le b$ gives $-a<b$. Together these conclusions are
$-a<b<a$.

\step{Part (a), Step 2: prove the reverse implication.}
Assume $-a<b<a$.
If $b\ge0$, then $\abs{b}=b<a$ by the right-hand assumed inequality.
If $b<0$, negate the left-hand assumed inequality $-a<b$ to get
$-b<a$; since $\abs{b}=-b$ in this case, we again obtain $\abs{b}<a$.
The case $b=0$ belongs to the first case.

\step{Part (a), Step 3: check nonpositive values of the bound.}
If $a\le0$, the left condition cannot hold: it would give
$0\le\abs{b}<a\le0$. The right condition cannot hold either.
Indeed $a\le0$ implies $0\le-a$, hence $a\le-a$, while
$-a<b<a$ would imply $-a<a$. Those inequalities contradict each other.
Thus the equivalence remains true when both conditions are false.

\answer{For all real $a,b$,
\(\boxed{\abs{b}<a\ \Longleftrightarrow\ -a<b<a}\).
Either condition, when true, forces $a>0$.}

\step{Part (b), Step 1: prove the forward implication.}
Assume $\abs{a-b}<c$. Apply part (a) to the real number $a-b$ with
bound $c$. It gives
\[
 -c<a-b<c.
\]
Add $b$ to the three expressions. Addition preserves strict inequality:
O4 preserves the weak inequality, and equality after addition would
contradict additive cancellation. Thus
\[
 -c+b<(a-b)+b<c+b.
\]
A1 and A4 give $(a-b)+b=a$; A2 gives $-c+b=b-c$ and $c+b=b+c$.
Therefore $b-c<a<b+c$.

\step{Part (b), Step 2: prove the reverse implication.}
Assume $b-c<a<b+c$. Add $-b$ to every expression:
\[
 (b-c)-b<a-b<(b+c)-b.
\]
Associativity and commutativity let $b$ and $-b$ become adjacent;
A4 and A3 then give $(b-c)-b=-c$ and $(b+c)-b=c$.
We obtain $-c<a-b<c$, and part (a) gives $\abs{a-b}<c$.

\answer{For all real $a,b,c$,
\(\boxed{\abs{a-b}<c\ \Longleftrightarrow\ b-c<a<b+c}\).
If $c\le0$, both conditions are false, as follows from part (a);
thus no unstated assumption $c>0$ is needed for the equivalence.}

\step{Part (c), Step 1: prove the forward implication.}
Assume $\abs{a-b}\le c$. Exercise~3.5(a), with bounded number $a-b$
and bound $c$, gives
\[
 -c\le a-b\le c.
\]
Use O4 to add $b$ throughout:
\[
 -c+b\le(a-b)+b\le c+b.
\]
The same algebra as in part (b) simplifies the three expressions,
giving $b-c\le a\le b+c$.

\step{Part (c), Step 2: prove the reverse implication.}
Assume $b-c\le a\le b+c$. Add $-b$ throughout by O4 and simplify:
\[
 (b-c)-b\le a-b\le(b+c)-b
 \quad\Longrightarrow\quad
 -c\le a-b\le c.
\]
Exercise~3.5(a) now gives $\abs{a-b}\le c$.

\step{Part (c), Step 3: check the boundary cases.}
If $c<0$, an absolute value cannot be at most $c$, so the left condition
is false. Also $b-c>b+c$, since subtracting gives
$(b-c)-(b+c)=-2c>0$; hence no $a$ can lie between those endpoints in
the specified order. The right condition is false as well.
If $c=0$, the left condition forces $a-b=0$, or $a=b$; the right
condition becomes $b\le a\le b$, which also forces $a=b$ by O2.

\answer{For all real $a,b,c$,
\(\boxed{\abs{a-b}\le c\ \Longleftrightarrow\ b-c\le a\le b+c}\).}

\exercise{3.8}{Comparison with every number above a fixed bound}

\goal Suppose that $a,b\in\R$ and
\[
 a\le b_1\qquad\text{for every real number $b_1$ satisfying $b_1>b$}.
 \tag{6}
\]
Prove $a\le b$.

\idea If $a$ were above $b$, the arithmetic midpoint between $a$ and $b$
would be above $b$ but still below $a$. That specific choice would
contradict (6). We will calculate both comparisons explicitly.

\step{Step 1: negate the desired conclusion.}
Suppose, for a contradiction, that $a\le b$ is false. By O1, either
$a\le b$ or $b\le a$. Since the first alternative is excluded,
$b\le a$. Equality $a=b$ would imply $a\le b$, so equality is also
excluded. Consequently $b<a$, or equivalently $a>b$.
Adding $-b$ to this strict inequality gives
\[
 d:=a-b>0.
\]

\step{Step 2: justify division by two and its sign.}
Exercise~3.4 gives $1>0$. Adding $1$ to $0<1$ gives $1<2$, and
transitivity with $0<1$ gives $2>0$. Thus $2\ne0$, so its inverse
$1/2$ exists and is positive by Theorem~3.2(vi).
Multiplying $0<d$ by the positive number $1/2$ preserves strict
inequality, by the argument in Exercise~3.4(vii), Step~2. Hence
\[
 \frac d2>0.
\]

\step{Step 3: construct a number to which the hypothesis applies.}
Define
\[
 m=\frac{a+b}{2}.
\]
We compute its difference from $b$ using a common denominator:
\begin{align*}
 m-b
   &=\frac{a+b}{2}-\frac{2b}{2}
       &&\text{because $2\ne0$ and $(2b)/2=b$},\\
   &=\frac{a+b-2b}{2}
       &&\text{by DL and the definition of division},\\
   &=\frac{a-b}{2}
       &&\text{since $b-2b=-b$},\\
   &=\frac d2>0
       &&\text{by the definition of $d$ and Step~2}.
\end{align*}
Adding $b$ to $m-b>0$ gives $m>b$. Thus $b_1=m$ satisfies the
restriction in (6).

\step{Step 4: show that this same number is below $a$.}
Again use a common denominator:
\begin{align*}
 a-m
   &=\frac{2a}{2}-\frac{a+b}{2},\\
   &=\frac{2a-a-b}{2},\\
   &=\frac{a-b}{2},\\
   &=\frac d2>0.
\end{align*}
Adding $m$ throughout $a-m>0$ gives $a>m$, or $m<a$.

\step{Step 5: apply the universal hypothesis and exhibit the contradiction.}
Since $m>b$, (6) applies with $b_1=m$ and gives $a\le m$.
Step~4 gave $m<a$, which includes $m\le a$ and $m\ne a$.
From $a\le m$ and $m\le a$, O2 gives $a=m$, contrary to $m\ne a$.
The assumption $a>b$ leads to a contradiction. Therefore $a\le b$.

\answer{The hypothesis forces \(\boxed{a\le b}\).
Only ordered-field arithmetic was used: the point $m$ was constructed
directly, without a completeness axiom or an appeal to a limit.}
